\section{Discussion and Conclusion}
\label{sec:limits}

The decomposition separates two benchmarks that a headline accuracy figure does not. On HC3, orthography alone matches content alone, 3.20\% against 3.26\% test error, and a fine-tuned DeBERTa improves on either by roughly 0.03 F1. On DAIGT V2, content is stronger than orthography by a factor of 7.9 in error, and a bag-of-words model over surface-normalised text comes within 0.0016 F1 of the best transformer, inside our measured seed range. A practitioner choosing an evaluation corpus, or a reviewer reading a reported score, cannot recover that distinction from the score itself.

What the measurement does not license is a claim that either corpus is clean. DAIGT V2's surface-only arm reaches 0.9214 F1, far above chance, so surface form is substantially informative on both. The finding is comparative, that surface reaches parity with content on one corpus and not the other, and it is bounded by the arms we built. A richer surface featurisation would raise both surface numbers, and a stronger content model would raise both content numbers, so the arms bound rather than partition the separability.

Two results here began as the opposite conclusion, and both are worth stating for what they say about method. The whitespace cue in HC3 is the most-discussed artefact in this literature, and it is sufficient in isolation. It is also unnecessary, since BERT's tokeniser cannot represent it and BERT reaches 0.9916 F1 regardless. A cleaning experiment that reports no change for BERT is therefore uninterpretable without the accompanying evidence that the pipeline runs at all, which we supply by showing the same pipeline does change BERT's DAIGT V2 predictions where the artefact removed is one WordPiece encodes. Similarly, a character perturbation drives DeBERTa from 0.9972 to 0.3644 F1 with a one-directional confusion matrix, which reads as a targeted vulnerability until the same model is shown to label uniform random strings as human in 100\% of cases at 0.98 confidence. We report the requirement for a label-free control rather than the vulnerability, and we suggest that comparable published collapses be read with that control in mind.

Several limits bound every number above. The transformers see 128 tokens, which covers a median HC3 document but only about 31\% of a median DAIGT V2 essay, so the DAIGT V2 transformer results describe classification from an essay's opening and the cross-corpus comparison is not like-for-like in the information each model receives. Seed spread is estimated from three seeds and reported as a range, not a test statistic, and one such estimate moved by 31\% when its runs were repeated, which is why no claim here rests on a variance ratio. HC3 contains output from one generator collected at one time, both corpora are English, and the classical arms share one preprocessing pipeline whose stopword removal and lemmatisation remove more than punctuation, so the content arm is a lower bound on content-carried separability rather than a clean isolate. We evaluate two corpora and five model families, and nothing here establishes that the pattern holds beyond them.

The measurement itself is the portable part. It requires no new annotation, runs in minutes on a laptop, and yields one number per corpus that a headline accuracy cannot express. Reporting it alongside detection results would let a reader see, before trusting a 99\% figure, how much of the corpus a model could pass without reading the language.
